\documentclass[10pt,a4paper]{amsart}
\usepackage[margin=19mm]{geometry}
\usepackage{amsmath,amssymb,mathtools}
\usepackage[scr=rsfs,cal=euler]{mathalfa}
\usepackage{xfrac}
\usepackage[unicode,hidelinks,bookmarks=false]{hyperref}
\newcommand{\ie}{i.e.\ }
\newcommand{\cf}{cf.\ }
\newcommand{\resp}{respectively\ }
\newcommand{\Subsection}{\subsection}
\providecommand{\nobreakdash}{\nobreak\hbox{-}}
\setlength{\emergencystretch}{2em}
\multlinegap0pt
\usepackage{amssymb}
\usepackage{xcolor}
\usepackage{algorithm}\usepackage{algpseudocode}
\usepackage{tikz}
\newtheorem{lemma}{Lemma}
\newtheorem{theorem}{Theorem}
\newcommand{\Adag}[1]{#1^{\dagger}}
\newcommand{\Perp}[1]{#1^{\perp}}
\newcommand{\trace}[1]{\textnormal{Tr}\left({#1}\right)}
\title{Information Geometry for Wasserstein KL Divergence of Gaussian Measures on \(\mathbb{R}^n\)}
\author{Adwait Datar, Nihat Ay}
\date{Source: arXiv:2503.24022v3 (CC BY 4.0)}
\begin{document}
\maketitle

\noindent\emph{Source and licence.} Information Geometry for Wasserstein KL Divergence of Gaussian Measures on \(\mathbb{R}^n\). Version arXiv:2503.24022v3 (CC BY 4.0), distributed by arXiv under the Creative Commons Attribution 4.0 International licence, http://creativecommons.org/licenses/by/4.0/, as recorded in the arXiv licence metadata for this exact version and shown on the abstract page https://arxiv.org/abs/2503.24022v3. Full licence notice: \texttt{licenses/arxiv-2503.24022-CC-BY-4.0.txt}.\par\smallskip

\noindent\textit{Editorial scope.} Counted unit: the article's theorem theom:WKL\_continuity (source0.tex lines 2562--2612). The article's complete original proof of it is delivered with it (source0.tex lines 2613--2714, one \texttt{proof} environment of 102 lines). No mathematics is written, repaired or re-derived: the statement and the proof are the article's own lines, in the article's own notation, and no retained line is substituted or re-flowed.  Retained uncounted as the source-local context the proof needs: lines 1473--1475; lines 2425--2427; lines 2541--2548.  Nothing was omitted from any retained span, so the package carries no omission record.  No retained line contains a citation, so the wrapper carries no bibliography and no bibliography entry.  No retained line contains a figure environment, an image-inclusion command or drawing code, so no image is delivered and the package declares no asset dependency.  Editorial formatting only, in addition: an AMS \texttt{amsart} preamble that restates the article's own declarations and macros used by the retained lines, namely the environments \texttt{lemma}, \texttt{theorem} on separate counters, together with the standard packages those lines need.  The retained environments print in this document as Lemma 1, Theorem 1 in order of appearance, whereas the article prints the same items with the numbering of its own full document.  The complete native source of the article as distributed in the arXiv e-print for this exact version is bundled unchanged as \texttt{sources/175687/source0.tex}.
\begin{align}\label{eq:D_main_formula}  
    D_{\rm WKL}(\mu \| \nu) &= \frac{1}{4}\trace{\Sigma_0-\Sigma_1+\Sigma_0R^2\log(R^2)} +\frac{1}{4}\left\lVert \sqrt{Q+2\Perp{\log(R)}} (m_1-m_0) \right\rVert^2
\end{align}

\begin{align}\label{eq:uniform_lower_bound}  
    D_{\rm WKL}\left(\mathcal{N}(m_0,\Sigma_0) \| \mathcal{N}(m_1,\Sigma_1)\right) &\geq  \frac{1}{4}\left\lVert m_1-m_0 \right\rVert^2
\end{align}

\begin{lemma}[Continuity of the spectral functional calculus]\label{lemm:continuous_spectral_calculus}
Let $f:[0,\infty)\to\mathbb R$ be continuous. For $X\succeq0$, define $f(X)$ by applying $f$ to the eigenvalues of $X$. Then the map
\[
\mathrm{cl}(S_n^+)\longrightarrow S_n,
\qquad X\longmapsto f(X),
\]
is continuous. In particular, $X\mapsto\chi(X)$ and $X\mapsto\kappa(X)$ are continuous on $\mathrm{cl}(S_n^+)$.
\end{lemma}

\begin{theorem}\label{theom:WKL_continuity}
Equip $\mathbb R^n$ and $S_n$ with their Euclidean topologies, all product spaces below with the product topology, and $[0,\infty]$ with its usual extended-real topology. Define
\[
\mathcal{L}(m_0,\Sigma_0,m_1,\Sigma_1)
=D_{\rm WKL}\!\left(\mathcal{N}(m_0,\Sigma_0)\,\middle\|\,\mathcal{N}(m_1,\Sigma_1)\right)
\]
on $\mathbb{R}^n\times S_n^+\times\mathbb{R}^n\times S_n^+$. Then the following statements hold.

\begin{enumerate}
\item $\mathcal{L}$ is continuous on its domain.

\item $\mathcal{L}$ has a unique continuous extended-real-valued extension $\overline{\mathcal{L}}$ to
\begin{align*}
\mathcal{U}
={}&\left(\mathbb{R}^n\times S_n^+\times\mathbb{R}^n\times\mathrm{cl}(S_n^+)\right)\\
&\quad\cup\left(\mathbb{R}^n\times\mathrm{cl}(S_n^+)\times\mathbb{R}^n\times S_n^+\right).
\end{align*}
If $\Sigma_0\succ0$ and $\Sigma_1\succeq0$, set
\[
R=\Sigma_0^{-1/2}\left(\Sigma_0^{1/2}\Sigma_1\Sigma_0^{1/2}\right)^{1/2}\Sigma_0^{-1/2}.
\]
Then
\begin{align}\label{eq:continuous_boundary_extension}
\overline{\mathcal{L}}(m_0,\Sigma_0,m_1,\Sigma_1)
={}&\frac14\trace{\Sigma_0-\Sigma_1+\Sigma_0\chi(R)}\nonumber\\
&+\frac14(m_1-m_0)^\top\kappa(R)(m_1-m_0).
\end{align}
If $\Sigma_0$ is singular and $\Sigma_1\succ0$, then $\overline{\mathcal{L}}(m_0,\Sigma_0,m_1,\Sigma_1)=+\infty$.

\item There is no continuous extension of $\overline{\mathcal{L}}$ to
\[
\mathcal{X}=\mathbb{R}^n\times\mathrm{cl}(S_n^+)\times\mathbb{R}^n\times\mathrm{cl}(S_n^+).
\]
One lower-semicontinuous extension to $\mathcal{X}$ is
\begin{align}\label{eq:lsc_extension}
\overline{\mathcal{L}}_{\rm ext}(m_0,\Sigma_0,m_1,\Sigma_1)
=\begin{cases}
\overline{\mathcal{L}}(m_0,\Sigma_0,m_1,\Sigma_1),
& \Sigma_0\succ0\ \textnormal{or}\ \Sigma_1\succ0,\\[5pt]
\dfrac14\lVert m_1-m_0\rVert^2,
& \Sigma_0,\Sigma_1\ \textnormal{both singular}.
\end{cases}
\end{align}
At a pair of Dirac measures, this chosen value equals the interior lower limit:
\begin{align}\label{eq:dirac_liminf}
\liminf_{\substack{\Sigma_0,\Sigma_1\succ0\\(\Sigma_0,\Sigma_1)\to(0,0)}}
\mathcal{L}(m_0,\Sigma_0,m_1,\Sigma_1)
=\frac14\lVert m_1-m_0\rVert^2.
\end{align}
\end{enumerate}
\end{theorem}

\begin{proof}
\noindent\textit{Step 1: scalar reduction and interior continuity.}
For $R\succ0$, diagonalizing $R$ shows that
\begin{align}\label{eq:kappa_identity}
\Adag{(R-I)}\left(\log(R^2)R^2-R^2+I\right)\Adag{(R-I)}
+2\Perp{\log R}
=\kappa(R).
\end{align}
Indeed, the scalar multiplier is $\kappa(r)$ on every eigenspace of $R$. The relevant scalar limits are
\[
\lim_{r\downarrow0}r^2\log(r^2)=0,
\qquad
\lim_{r\downarrow0}\kappa(r)=1,
\qquad
\lim_{r\to1}\kappa(r)=2.
\]
Thus $\chi$ and $\kappa$ are continuous on $[0,\infty)$. Lemma~\ref{lemm:continuous_spectral_calculus} implies that $X\mapsto\chi(X)$ and $X\mapsto\kappa(X)$ are continuous on the cone of positive-semidefinite matrices.

\medskip
\noindent\textit{Step 2: target-covariance boundary.}
When $\Sigma_0\succ0$, the map
\[
(\Sigma_0,\Sigma_1)\longmapsto
\Sigma_0^{-1/2}\left(\Sigma_0^{1/2}\Sigma_1\Sigma_0^{1/2}\right)^{1/2}\Sigma_0^{-1/2}
\]
is continuous for $\Sigma_1\succeq0$. Hence \eqref{eq:continuous_boundary_extension} is finite and continuous when the first covariance is positive definite, including when $\Sigma_1$ is singular. On the interior, \eqref{eq:kappa_identity} shows that this formula agrees with \eqref{eq:D_main_formula}. This proves the first assertion and the target-covariance part of the second assertion.

\medskip
\noindent\textit{Step 3: source-covariance boundary.}
Consider a sequence
\[
\Sigma_{0,k}\to\overline{\Sigma}_0,
\qquad
\Sigma_{1,k}\to\overline{\Sigma}_1,
\]
where $\overline{\Sigma}_0$ is singular and $\overline{\Sigma}_1\succ0$. Let $R_k$ be the corresponding matrix. For all sufficiently large $k$, there is a constant $c>0$ such that $\Sigma_{1,k}\succeq cI$. Let $v_k$ be a unit eigenvector of $\Sigma_{0,k}$ associated with $\lambda_{\min}(\Sigma_{0,k})$. By operator monotonicity of the positive-semidefinite square root,
\[
\left(\Sigma_{0,k}^{1/2}\Sigma_{1,k}\Sigma_{0,k}^{1/2}\right)^{1/2}
\succeq\sqrt{c}\,\Sigma_{0,k}^{1/2}.
\]
It follows that
\[
v_k^\top R_kv_k
\geq\sqrt{\frac{c}{\lambda_{\min}(\Sigma_{0,k})}}
\longrightarrow+\infty.
\]
Consequently, if $r_k=\lambda_{\max}(R_k)$, then $r_k\to+\infty$.

Set $M_k=\log(R_k^2)R_k^2-R_k^2+I$. Since $R_k\Sigma_{0,k}R_k=\Sigma_{1,k}$, an orthonormal eigenbasis $\{u_{j,k}\}$ of $R_k$, with eigenvalues $r_{j,k}$, gives
\begin{align*}
\trace{M_k\Sigma_{0,k}}
&=\sum_j\left(2\log r_{j,k}-1+r_{j,k}^{-2}\right)u_{j,k}^\top\Sigma_{1,k}u_{j,k}.
\end{align*}
Every summand is nonnegative because $z-1-\log z\geq0$ for $z>0$, applied with $z=r_{j,k}^{-2}$. Keeping only an eigenvector associated with $r_k$ therefore yields
\[
\trace{M_k\Sigma_{0,k}}
\geq c\left(2\log r_k-1+r_k^{-2}\right)
\longrightarrow+\infty.
\]
The covariance contribution alone therefore diverges, proving continuity with value $+\infty$ on the source-singular part of $\mathcal{U}$. Uniqueness follows because the original domain is dense in $\mathcal{U}$ and $[0,\infty]$ is Hausdorff.

\medskip
\noindent\textit{Step 4: obstruction to a full continuous extension.}
Write $\Delta m=m_1-m_0$. Along $\Sigma_0=\Sigma_1=\varepsilon I$,
\[
\mathcal{L}=\frac12\lVert\Delta m\rVert^2.
\]
Along $\Sigma_0=\varepsilon I$ and $\Sigma_1=\varepsilon^2I$,
\[
\mathcal{L}
=\frac{n}{4}\left(\varepsilon-\varepsilon^2+\varepsilon^2\log\varepsilon\right)
+\frac14\kappa(\sqrt{\varepsilon})\lVert\Delta m\rVert^2
\longrightarrow\frac14\lVert\Delta m\rVert^2.
\]
These paths have different limits when $\Delta m\neq0$. If $\Delta m=0$, compare the diagonal path with
\[
\Sigma_0=\varepsilon e^{-1/\varepsilon}I,
\qquad
\Sigma_1=\varepsilon I.
\]
Along the latter path,
\[
\mathcal{L}
=\frac{n}{4}\left(\varepsilon e^{-1/\varepsilon}-\varepsilon+1\right)
\longrightarrow\frac n4,
\]
whereas the diagonal limit is zero. Thus continuity fails at every point $(m_0,0,m_1,0)$.

\medskip
\noindent\textit{Step 5: lower-semicontinuous extension and the Dirac lower limit.}
Since $\mathcal{U}$ is open in $\mathcal{X}$ and $\overline{\mathcal{L}}$ is continuous on $\mathcal{U}$, it remains only to check lower semicontinuity at points where both covariance matrices are singular. The uniform bound \eqref{eq:uniform_lower_bound} passes to the one-sided boundary extension:
\[
\overline{\mathcal{L}}(m_0,\Sigma_0,m_1,\Sigma_1)
\geq\frac14\lVert m_1-m_0\rVert^2
\]
whenever $(m_0,\Sigma_0,m_1,\Sigma_1)\in\mathcal{U}$; on $\mathcal X\setminus\mathcal U$, the same bound holds with equality by \eqref{eq:lsc_extension}. If a sequence converges to a point at which both covariance matrices are singular, this bound and the continuity of the means imply
\[
\liminf_k\overline{\mathcal{L}}_{\rm ext}(m_{0,k},\Sigma_{0,k},m_{1,k},\Sigma_{1,k})
\geq\frac14\lVert m_1-m_0\rVert^2.
\]
This is precisely the value assigned in \eqref{eq:lsc_extension}; hence the extension is lower semicontinuous. Equation \eqref{eq:dirac_liminf} follows from the uniform lower bound and the path $\Sigma_0=\varepsilon I$, $\Sigma_1=\varepsilon^2I$.
\end{proof}

\end{document}
